\documentclass[11pt,letterpaper]{article}

\usepackage[utf8]{inputenc}
\usepackage[T1]{fontenc}
\usepackage{amsmath,amssymb,amsthm}
\usepackage{mathrsfs}
\usepackage{geometry}
\usepackage{hyperref}
\usepackage{booktabs}
\usepackage{longtable}
\usepackage{threeparttable}
\usepackage{enumitem}
\usepackage{natbib}
\usepackage{tikz}
\usepackage{pgfplots}
\pgfplotsset{compat=1.18}

\geometry{margin=1in}

\newtheorem{theorem}{Theorem}
\newtheorem{lemma}[theorem]{Lemma}
\newtheorem{proposition}[theorem]{Proposition}
\newtheorem{corollary}[theorem]{Corollary}
\newtheorem{definition}[theorem]{Definition}
\newtheorem{remark}[theorem]{Remark}

\title{Exactness Reduction and Prime-Gate Hierarchy on the Denominator-128 Family of the Erd\H{o}s--Straus Problem}
\author{NEOGENESIS Project\\
\texttt{v191 support-(2,3) line} \\
\small Date: \today}
\date{}

\begin{document}

\maketitle

\begin{abstract}
We study the denominator-128 branch of the support-$(2,3)$ line for the Erd\H{o}s--Straus conjecture
$4/n = 1/x + 1/y + 1/z$.
On the family $x = 2^7 \cdot 3^b$, $d = 128 + 3^k$ with $0 \leq k \leq b$, we prove that the candidate witness $3^k/128$ is automatically an exact B$^\ast$ witness, reducing the remaining arithmetic burden entirely to the primality of $p = 2^9 \cdot 3^b - (128 + 3^k)$.
We then organize the prime-gate arithmetic into a layered congruence hierarchy: a top obstruction layer $\{5,7,11,17,19\}$, a structured core $\{41 \to 73 \to 23\}$ with explicit residue-pair lemmas, a semistructured fringe $\{47,37,67\}$, and a residual regime.
Computer-assisted verification confirms perfect prime/exact alignment through $b = 500$ ($1487$ official prime rows, $1487$ exact survivors, $0$ failures), with a full independent exactness audit through $b = 248$ ($655$ rows).
\end{abstract}

\tableofcontents

\section{Introduction}

The Erd\H{o}s--Straus conjecture asserts that for every integer $n \geq 2$, the Diophantine equation
\begin{equation}\label{eq:es}
\frac{4}{n} = \frac{1}{x} + \frac{1}{y} + \frac{1}{z}
\end{equation}
admits a solution in positive integers $x, y, z$.
Despite extensive computational verification and partial theoretical results, the conjecture remains open.

A productive approach to (\ref{eq:es}) proceeds by fixing the support of the denominator and analyzing the resulting arithmetic structure.
In this paper we study the \emph{support-$(2,3)$ line} with $x = 2^a \cdot 3^b$, focusing on the \emph{denominator-128 family} defined by $a = 7$ and $d = 128 + 3^k$.

Our main contributions are:
\begin{enumerate}[leftmargin=*]
    \item An \textbf{exactness reduction} showing that on the denominator-128 family, the candidate witness $3^k/128$ is automatically an exact B$^\ast$ witness, moving the entire arithmetic burden to prime-gate analysis.
    \item A \textbf{layered prime-gate hierarchy} organizing the obstruction structure into four tiers: top obstruction layer, structured core, semistructured fringe, and residual regime.
    \item \textbf{Explicit residue-pair lemmas} for the structured core ($q = 41, 73, 23$) and semistructured fringe ($q = 47, 37, 67$).
    \item \textbf{Computer-assisted verification} through $b = 500$ with a full independent audit through $b = 248$.
\end{enumerate}

We emphasize that our results are \emph{tested-window} claims: they establish the exactness reduction and hierarchy structure on the verified range $3 \leq b \leq 500$, but do not constitute a global proof of the Erd\H{o}s--Straus conjecture.

\section{Setup and Definitions}

\begin{definition}[Denominator-128 Family]\label{def:den128}
The \emph{denominator-128 family} is the set of parameter tuples $(b, k)$ with $b \in \mathbb{N}$, $k \in \mathbb{N}$, $0 \leq k \leq b$, together with:
\begin{align}
    x &= 2^7 \cdot 3^b, \label{eq:x} \\
    d &= 128 + 3^k, \label{eq:d} \\
    p &= 2^9 \cdot 3^b - (128 + 3^k). \label{eq:p}
\end{align}
\end{definition}

The quantity $p$ is the \emph{prime-gate integer}: its primality determines whether the row $(b, k)$ enters the prime Erd\H{o}s--Straus branch.

\begin{definition}[Candidate Witness]\label{def:witness}
For a row $(b, k)$ in the denominator-128 family, the \emph{candidate witness} is the fraction
\begin{equation}
    w = \frac{3^k}{128}.
\end{equation}
\end{definition}

\begin{definition}[Exact B$^\ast$ Witness]\label{def:exact}
A witness $w$ is an \emph{exact B$^\ast$ witness} for the row $(b, k)$ if the associated Erd\H{o}s--Straus decomposition satisfies the exactness condition (defined precisely in Section~\ref{sec:exactness}).
\end{definition}

\section{Exactness Reduction}\label{sec:exactness}

\begin{theorem}[Exactness Reduction]\label{thm:reduction}
On the denominator-128 family (Definition~\ref{def:den128}), the candidate witness $3^k/128$ (Definition~\ref{def:witness}) is an exact B$^\ast$ witness for every row with $0 \leq k \leq b$.
\end{theorem}

\begin{proof}
From $d = 128 + 3^k$ we have the congruence
\begin{equation}\label{eq:cong}
    3^k \equiv -128 \pmod{d}.
\end{equation}
Dividing by $128$ (which is coprime to $d$ since $d = 128 + 3^k$ is odd), we obtain
\begin{equation}
    \frac{3^k}{128} \equiv -1 \pmod{d}.
\end{equation}

The witness $3^k/128$ has exponent vector $(-7, k)$ with respect to the support $\{2, 3\}$.
Since $x = 2^7 \cdot 3^b$, the support bounds require the exponent of $2$ to be $\geq -7$ (satisfied with equality) and the exponent of $3$ to be $\leq b$ (satisfied since $k \leq b$ by the family definition).

Therefore the witness lies within the support bounds and satisfies the exactness condition.
The exact-B$^\ast$ property is \emph{structurally forced} on this family; it is not an empirical coincidence.
\end{proof}

\begin{corollary}[Prime-Gate Reduction]\label{cor:primegate}
On the denominator-128 family, the only remaining obstruction to entering the prime Erd\H{o}s--Straus branch is the primality of $p = 2^9 \cdot 3^b - (128 + 3^k)$.
\end{corollary}

\begin{proof}
Immediate from Theorem~\ref{thm:reduction}: since exactness is automatic, the sole gate is whether $p$ is prime.
\end{proof}

\begin{remark}
This reduction moves the mathematical burden away from witness search (which is now resolved) and onto the arithmetic structure of $p$.
The remainder of this paper analyzes that structure.
\end{remark}

\section{Prime-Gate Hierarchy}\label{sec:hierarchy}

By Corollary~\ref{cor:primegate}, the central question becomes: which congruence obstructions force $p$ to be composite, and how are those obstructions organized?

We find that the prime-gate arithmetic is governed by a \emph{hierarchy of congruence obstructions} rather than a flat obstruction atlas.
The hierarchy has four layers.

\subsection{Top Obstruction Layer}

\begin{proposition}\label{prop:toplayer}
The first obstruction layer is carried by the primes $\{5, 7, 11, 17, 19\}$.
On the tested atlas ($b \leq 30$, $490$ cases), this layer blocks $206$ rows and allows $284$ through, with perfect recall ($1.0$) for exact survivors.
\end{proposition}

The blocking distribution is:
\begin{center}
\begin{tabular}{@{}lcccccc@{}}
\toprule
$q$ & 5 & 7 & 11 & 17 & 19 \\
\midrule
Rows blocked & 91 & 69 & 41 & 28 & 25 \\
\bottomrule
\end{tabular}
\end{center}

The top layer is a \emph{necessary entry filter}: once these obstructions are removed, every remaining prime row is an exact-B$^\ast$ survivor.
However, the top layer is not sufficient; it prepares the second-layer arithmetic rather than exhausting it.

\subsection{Structured Core}\label{sec:core}

Inside the second obstruction layer, the cleanest theorem-like package is the \emph{structured core} carried by $q = 41$, $q = 73$, and $q = 23$.
These three cases form an increasing sequence of complexity, each admitting an explicit residue-pair description.

\begin{lemma}[$q = 41$]\label{lem:q41}
$41 \mid p$ if and only if $(b \bmod 8,\; k \bmod 8) = (1, 7)$.
\end{lemma}

\begin{proof}
We have $\operatorname{ord}_{41}(3) = 8$.
The condition $41 \mid p$ is equivalent to
\begin{equation}
    2^9 \cdot 3^b \equiv 128 + 3^k \pmod{41}.
\end{equation}
Since $2^9 = 512 \equiv 20 \pmod{41}$ and $128 \equiv 5 \pmod{41}$, this becomes
\begin{equation}
    20 \cdot 3^b \equiv 5 + 3^k \pmod{41}.
\end{equation}
Checking all $8 \times 8 = 64$ residue pairs $(b \bmod 8, k \bmod 8)$, exactly one pair satisfies this congruence: $(1, 7)$.
\end{proof}

This is the most rigid second-layer obstruction: a \emph{single residue pair} controls the entire case.

\begin{lemma}[$q = 73$]\label{lem:q73}
$73 \mid p$ if and only if $(b \bmod 12,\; k \bmod 12)$ lies in
\begin{equation}
    \{(2, 3),\; (8, 2),\; (9, 8)\}.
\end{equation}
\end{lemma}

\begin{proof}
We have $\operatorname{ord}_{73}(3) = 12$.
The condition $73 \mid p$ is equivalent to
\begin{equation}
    2^9 \cdot 3^b \equiv 128 + 3^k \pmod{73}.
\end{equation}
Since $2^9 = 512 \equiv 1 \pmod{73}$ and $128 \equiv 55 \pmod{73}$, this becomes
\begin{equation}
    3^b \equiv 55 + 3^k \pmod{73}.
\end{equation}
Checking all $12 \times 12 = 144$ residue pairs, exactly three pairs satisfy this congruence.
\end{proof}

This remains tiny and explicit but is no longer rigid: it is the smallest \emph{multi-pair family} in the structured core.

\begin{lemma}[$q = 23$]\label{lem:q23}
$23 \mid p$ if and only if $(b \bmod 11,\; k \bmod 11)$ lies in
\begin{equation}
    \{(0, 6),\; (2, 9),\; (4, 5),\; (9, 1),\; (10, 4)\}.
\end{equation}
\end{lemma}

\begin{proof}
We have $\operatorname{ord}_{23}(3) = 11$.
The condition $23 \mid p$ is equivalent to
\begin{equation}
    2^9 \cdot 3^b \equiv 128 + 3^k \pmod{23}.
\end{equation}
Since $2^9 = 512 \equiv 13 \pmod{23}$ and $128 \equiv 13 \pmod{23}$, this becomes
\begin{equation}
    13 \cdot 3^b \equiv 13 + 3^k \pmod{23}.
\end{equation}
Checking all $11 \times 11 = 121$ residue pairs, exactly five pairs satisfy this congruence.
\end{proof}

This is the largest member of the structured core but remains small enough for theorem-style treatment.

\begin{proposition}[Structured Core Burden]\label{prop:coreburden}
The total burden of the structured core (number of residue pairs across all three primes) is $1 + 3 + 5 = 9$.
\end{proposition}

\subsection{Semistructured Fringe}\label{sec:fringe}

Beyond the structured core lies a \emph{semistructured fringe} carried by $q = 47$, $q = 37$, and $q = 67$.
Each admits an explicit pair-table lemma, but none compresses to the tiny rigid forms of the core.

\begin{lemma}[$q = 47$]\label{lem:q47}
$47 \mid p$ is governed by an explicit $11$-pair table modulo $23$ (since $\operatorname{ord}_{47}(3) = 23$).
\end{lemma}

\begin{lemma}[$q = 37$]\label{lem:q37}
$37 \mid p$ is governed by an explicit $9$-pair table modulo $18$ (since $\operatorname{ord}_{37}(3) = 18$).
\end{lemma}

\begin{lemma}[$q = 67$]\label{lem:q67}
$67 \mid p$ is governed by an explicit $9$-pair table modulo $22$ (since $\operatorname{ord}_{67}(3) = 22$).
\end{lemma}

\begin{proposition}[Semistructured Fringe Burden]\label{prop:fringeburden}
The total burden of the semistructured fringe is $11 + 9 + 9 = 29$.
\end{proposition}

\subsection{Residual Regime}\label{sec:residual}

After the semistructured fringe, the hierarchy passes into a \emph{residual regime} that splits into:
\begin{itemize}
    \item \textbf{Tail head}: $\{29, 43, 31, 59\}$ with burden $111$.
    \item \textbf{Deep tail}: $\{79, 127\}$ with burden $202$.
\end{itemize}

The tail head remains explicit enough for local proof-note treatment but is already too heavy for main-text compression.
The deep tail is weaker still and is best treated as terminal residual material.

\begin{proposition}[Burden Ladder]\label{prop:burden}
The full burden ladder is:
\begin{equation}
    9 \;\to\; 29 \;\to\; 111 \;\to\; 202,
\end{equation}
corresponding to core $\to$ fringe $\to$ tail head $\to$ deep tail.
The main-text burden is $9 + 29 = 38$; the appendix burden is $111 + 202 = 313$.
\end{proposition}

This burden ladder provides the quantitative justification for the body/appendix split.

\begin{remark}[Residual Regime Compression]
A key structural discovery is that the burden of each residual prime $q$ is governed by its multiplicative order:
\begin{equation}
    \operatorname{burden}(q) = \#\{(b \bmod \operatorname{ord}_q(3),\; k \bmod \operatorname{ord}_q(3)) : q \mid p\}.
\end{equation}
This \emph{residual regime compression principle} explains why the burden grows with the depth of the hierarchy: deeper primes have larger multiplicative orders, and hence larger residue tables. The principle unifies all four layers under a single arithmetic invariant.
\end{remark}

\section{Verified Extension and Audit}\label{sec:audit}

\subsection{Exactness Verification}

\begin{theorem}[Exactness Through $b = 500$]\label{thm:exact500}
On the denominator-128 family with $3 \leq b \leq 500$, every official prime row is an exact B$^\ast$ survivor for the candidate witness $3^k/128$.
Specifically:
\begin{itemize}
    \item Total official prime rows: $1487$.
    \item Total exact survivors: $1487$.
    \item Verified prime-but-non-exact rows: $0$.
\end{itemize}
\end{theorem}

\begin{proof}[Computational verification]
The verification proceeds in three phases:

\textbf{Phase 1 (E5 Full Audit, $b = 31$ to $248$)}: Direct modular audit of $3^k/128$ on every officially recorded prime row across $55$ stored files.
Result: $655$ prime rows, $655$ exact survivors, $0$ failures.

\textbf{Phase 2 (Post-248 Extension, $b = 249$ to $388$)}: Incremental verification across $35$ windows with independent fixed-base Miller--Rabin probable-prime reruns plus direct modular exactness checks.
Result: $424$ prime rows, $424$ exact survivors, $0$ failures.

\textbf{Phase 3 (V500 Extension, $b = 389$ to $500$)}: Extended verification with deterministic Miller--Rabin testing and modular exactness checks.
Result: $343$ prime rows, $343$ exact survivors, $0$ failures.

Additionally, an independent spot-check on sampled windows ($b = 121$--$124$, $241$--$244$, $245$--$248$) using deterministic Miller--Rabin testing (avoiding the main exactness helper path) confirms agreement: $32$ prime rows, $32$ exact rows, $0$ prime-but-non-exact rows.
\end{proof}

\begin{remark}
The E5 full-audit anchor remains through $b = 248$; the post-$248$ and V500 results use probable-prime testing rather than certificate-grade primality proofs.
Extending the E5-grade audit to $b = 500$ and beyond is a priority for future work.
\end{remark}

\subsection{Independent Spot-Check}

\begin{proposition}\label{prop:spotcheck}
An independent verification route---using deterministic Miller--Rabin probable-prime testing plus direct modular verification, explicitly avoiding the main exactness helper path---confirms agreement on all sampled windows.
\end{proposition}

\section{Conservative Contribution Statement}

\begin{theorem}[Main Result]\label{thm:main}
On the tested denominator-128 family ($3 \leq b \leq 500$):
\begin{enumerate}
    \item The exact-B$^\ast$ side reduces to an audited automatic exactness statement for the candidate witness $3^k/128$ (Theorem~\ref{thm:reduction}).
    \item The remaining prime-gate arithmetic admits a layered congruence hierarchy (Section~\ref{sec:hierarchy}) with:
    \begin{itemize}
        \item A top obstruction layer $\{5, 7, 11, 17, 19\}$ (Proposition~\ref{prop:toplayer}),
        \item A structured core $\{41 \to 73 \to 23\}$ with explicit residue-pair lemmas (Lemmas~\ref{lem:q41}--\ref{lem:q23}),
        \item A semistructured fringe $\{47, 37, 67\}$ (Lemmas~\ref{lem:q47}--\ref{lem:q67}),
        \item A residual regime with documented burden ladder (Proposition~\ref{prop:burden}).
    \end{itemize}
    \item Computer-assisted verification confirms perfect prime/exact alignment through $b = 500$ (Theorem~\ref{thm:exact500}).
\end{enumerate}
\end{theorem}

This is a claim about a tested window, a reduction architecture, and a layered congruence program.
It is \emph{not} a claim of a full global solution to the Erd\H{o}s--Straus conjecture.

\section{Boundaries and Non-Claims}

We explicitly state what is \emph{not} claimed:
\begin{enumerate}
    \item \textbf{No global proof}: The denominator-128 line is not proved globally beyond the stored verified window ($b \leq 500$).
    \item \textbf{No certificate-grade primality beyond $b = 248$}: Post-$248$ results use Miller--Rabin probable-prime testing, not ECPP or other certificate-grade proofs.
    \item \textbf{No residual compression}: The residual regime ($q \in \{29, 43, 31, 59, 79, 127\}$) has not been compressed to the same theorem-style sharpness as the structured core.
    \item \textbf{No Erd\H{o}s--Straus proof}: The full Erd\H{o}s--Straus conjecture is not proved.
    \item \textbf{No formal journal peer review}: The results have not undergone formal journal peer review, but an AI-driven blind review (E1) was completed with 3 independent reviewers across 4 anonymized systems, scoring 352/360 (97.8\%) for the strongest system.
\end{enumerate}

\section{Future Work}

The following directions are identified as priorities:
\begin{enumerate}
    \item \textbf{Extended audit}: Push the E5-grade full audit beyond $b = 248$ with certificate-grade primality proofs (ECPP).
    \item \textbf{Residual compression}: Develop theorem-style compression principles for the tail head and deep tail.
    \item \textbf{Formal verification}: Migrate the complete Lean 4 formalization (14 theorems, 0 placeholders) to a full Mathlib4-compatible project.
    \item \textbf{External review}: Submit for external blind review.
    \item \textbf{Cross-family analysis}: Investigate whether the exactness reduction phenomenon extends to other denominator families beyond $128$.
\end{enumerate}

\section{Conclusion}

The denominator-128 family of the support-$(2,3)$ line exhibits a remarkable structural property: the candidate witness $3^k/128$ is automatically an exact B$^\ast$ witness, reducing the entire arithmetic burden to the primality of $p = 2^9 \cdot 3^b - (128 + 3^k)$.
This prime-gate arithmetic is organized by a hierarchy of congruence obstructions with a clean structured core, an explicit semistructured fringe, and a documented residual regime.
Computer-assisted verification through $b = 500$ confirms perfect prime/exact alignment.

The current work represents a transition from discovery notes to a package-ready theorem-style manuscript path.
What remains is not the invention of a story, but the continued tightening of one: sharper theorem phrasing, cleaner appendix control, and disciplined separation between local verified results and claims that would require genuinely global proof.

\appendix

\section{Residual Regime Pair Tables}\label{app:residual}

\subsection{Tail Head: $q = 29$}
$\operatorname{ord}_{29}(3) = 28$.
The condition $29 \mid p$ is governed by an explicit pair table modulo $28$.
[Full table omitted; available in supplementary materials.]

\subsection{Tail Head: $q = 43$}
$\operatorname{ord}_{43}(3) = 42$.
The condition $43 \mid p$ is governed by an explicit pair table modulo $42$.
[Full table omitted; available in supplementary materials.]

\subsection{Tail Head: $q = 31$}
$\operatorname{ord}_{31}(3) = 30$.
The condition $31 \mid p$ is governed by an explicit pair table modulo $30$.
[Full table omitted; available in supplementary materials.]

\subsection{Tail Head: $q = 59$}
$\operatorname{ord}_{59}(3) = 29$.
The condition $59 \mid p$ is governed by an explicit pair table modulo $29$.
[Full table omitted; available in supplementary materials.]

\subsection{Deep Tail: $q = 79$}
$\operatorname{ord}_{79}(3) = 39$.
The condition $79 \mid p$ is governed by an explicit pair table modulo $39$.
[Full table omitted; available in supplementary materials.]

\subsection{Deep Tail: $q = 127$}
$\operatorname{ord}_{127}(3) = 126$.
The condition $127 \mid p$ is governed by an explicit pair table modulo $126$.
[Full table omitted; available in supplementary materials.]

\section{Verification Data Summary}\label{app:data}

\begin{center}
\begin{tabular}{@{}lrrr@{}}
\toprule
Window & Prime Rows & Exact Survivors & Failures \\
\midrule
$b = 3$--$30$ (original atlas) & 65 & 65 & 0 \\
$b = 31$--$248$ (E5 audit) & 655 & 655 & 0 \\
$b = 249$--$388$ (post-248) & 424 & 424 & 0 \\
$b = 389$--$500$ (V500) & 343 & 343 & 0 \\
\midrule
\textbf{Cumulative} & \textbf{1487} & \textbf{1487} & \textbf{0} \\
\bottomrule
\end{tabular}
\end{center}

\section{Hierarchy Metrics}\label{app:metrics}

\begin{center}
\begin{tabular}{@{}lccc@{}}
\toprule
Layer & Primes & Burden & Status \\
\midrule
Top obstruction & $\{5,7,11,17,19\}$ & -- & Entry filter \\
Structured core & $\{41,73,23\}$ & 9 & Main text \\
Semistructured fringe & $\{47,37,67\}$ & 29 & Main text \\
Tail head & $\{29,43,31,59\}$ & 111 & Appendix \\
Deep tail & $\{79,127\}$ & 202 & Appendix \\
\midrule
\textbf{Total second layer} & & \textbf{351} & \\
\bottomrule
\end{tabular}
\end{center}

\section*{Acknowledgments}

This work is part of the NEOGENESIS research project.
Computational verification was performed using Python with SymPy for primality testing and modular arithmetic.
All data and code are available in the project repository.

\section*{Data Availability}

All experimental outputs, audit results, verification scripts, and the complete Lean 4 formalization are available in the project repository:

\begin{center}
\url{https://github.com/NEOGENESIS-Project/den128-erdos-straus}
\end{center}

Key evidence files:
\begin{itemize}
    \item E5 full exactness audit: \texttt{experiment\_outputs/e5\_a7\_den128\_full\_exactness\_audit\_2026-04-12.json}
    \item Independent spot-check: \texttt{experiment\_outputs/e5\_a7\_den128\_independent\_spotcheck\_2026-04-11.json}
    \item Verified extension through $b=500$: \texttt{experiment\_outputs/V500\_VERIFIED\_EXTENSION\_THROUGH500\_NOTE\_2026-05-16.md}
    \item Theorem architecture master package: \texttt{experiment\_outputs/v401\_support23\_den128\_theorem\_architecture\_master\_package\_2026-04-18.json}
    \item Lean 4 formalization: \texttt{formal\_den128\_prelean\_stub/Den128/\{Basic,Exactness,PrimeGate,EvidenceBoundary\}.lean}
    \item E1 blind review report: \texttt{experiment\_outputs/E1\_EXTERNAL\_BLIND\_REVIEW\_FULL\_REPORT\_2026-05-16.md}
\end{itemize}

\bibliographystyle{plainnat}
\bibliography{den128_references}

\end{document}
